%! AP Statistics — Unit 2, Lesson 2.8 — Warm-Up
\documentclass[10pt]{article}
\usepackage{apstats-article}
\usepackage{apstats-boxes}

\begin{document}

\pageheader{Unit 2, Lesson 2.8}{Warm-Up}
\namedateperiod

\vspace{0.3in}

% ── PART 1 ────────────────────────────────────────────────────────────────────
{\large\bfseries\color{navy} Part 1 \quad Residuals (Lesson 2.7 Review)}

\vspace{0.12in}
\noindent Recall: a \textbf{residual} is the difference between an observed value and the
predicted value: $e = y - \hat{y}$.

\vspace{0.1in}

\begin{enumerate}[leftmargin=1.6em, itemsep=0.22in]

  \item A researcher models daily high temperature ($x$, \textdegree F) vs.\ ice cream sales
        ($\hat{y}$, dollars) with the equation $\widehat{\text{sales}} = -150 + 8(\text{temp})$.

        \begin{enumerate}[label=\textbf{(\alph*)}, leftmargin=1.6em, itemsep=8pt]
          \item Predict sales on a day when the temperature is 85\textdegree F.\\[8pt]
                \writeline

          \item On that day, actual sales were \$590. Calculate the residual and interpret it.\\[8pt]
                \writeline\\[1pt]\writeline

          \item The sum of all residuals for a least-squares line equals \blank{2cm}.
                Why does this make sense geometrically?\\[8pt]
                \writeline\\[1pt]\writeline
        \end{enumerate}

  \item Why do we minimize $\sum e_i^2$ instead of $\sum e_i$ when finding the best-fit line?\\[8pt]
        \writeline\\[1pt]\writeline

\end{enumerate}

\vspace{0.22in}

% ── PART 2 ────────────────────────────────────────────────────────────────────
{\large\bfseries\color{navy} Part 2 \quad Correlation Review (Lesson 2.5)}

\vspace{0.12in}

\begin{enumerate}[leftmargin=1.6em, itemsep=0.18in, start=3]

  \item A dataset has $r = 0.60$, $s_x = 5$, $s_y = 10$.

        \begin{enumerate}[label=\textbf{(\alph*)}, leftmargin=1.6em, itemsep=6pt]
          \item What is the value of $r^2$? \blank{3cm} What does this number represent (informally)?\\[8pt]
                \writeline\\[1pt]\writeline

          \item If the slope of the LSRL is computed as $b = r\,\dfrac{s_y}{s_x}$, what is $b$?
                Show the substitution.\\[8pt]
                \writeline\\[1pt]\writeline
        \end{enumerate}

\end{enumerate}

\vspace{0.14in}

\begin{tcolorbox}[colback=greenbg, colframe=greenacc, boxrule=0.8pt, arc=2mm,
    left=4mm, right=4mm, top=3mm, bottom=3mm]
\small\textbf{Getting ready for today:} The formula $b = r\,\dfrac{s_y}{s_x}$ is the
\emph{exact} slope of the least-squares line. Today you'll see \emph{why} this is the right
formula --- and what it means to ``minimize'' the sum of squared residuals.
\end{tcolorbox}

\end{document}
